\bookchapter{第四十五章：品質保証——AIが書いたコードを信頼できるものに}

\chapterintro{AIが書くコードにはひとつ特徴があります。たいてい「正しそうに見える」のです。変数の命名は規範的、コメントは揃っていて、ロジックも明快。しかし「正しそうに見える」は「本当に正しい」と同じではありません。AIコードに潜む問題は、手書きのコードより見つけにくいことがよくあります。表面の品質が高すぎて、人は油断しやすいからです。この章では、AIコードのために効果的な品質保証の仕組みを作る方法を教えます。テスト·レビュー·バージョン管理という3つの防線で、「正しそうに見える」を「本当に正しい」に変えましょう。}

\booksection{一、テスト駆動：AIにテストを書かせる}

テストはAIプログラミングの品質保証における第一の防線です。朗報があります。AIにテストを書かせるのは、機能のコードを書かせるより簡単で、信頼性も高いということです。テストの期待動作は明確なので、「正しそうに見えて実際は違う」コードが生まれにくいのです。

テスト駆動のAIプログラミングの流れはこうです。まずAIに機能のコードを書かせ、基本動作の検証が通ったら、すぐにテストケースを書かせます。テストケースは正常系、境界値、異常系をカバーすべきです。そのうえで、テストの結果に応じてAIに機能コードを修正させます。この流れの良いところは、テストが一度できあがれば、それが持続的な検証ツールになることです。その後のどんな変更も、テストを実行してリグレッションが起きていないか確認できます。

\begin{lstlisting}
calculate_discount関数のユニットテストを書いてください。
テストケースは以下をカバーすべきです：
1. 正常系：
   - 100元の商品を20%引き -> 80
   - 200元の商品を50%引き -> 100
2. 境界値：
   - 0元の商品を20%引き -> 0
   - 100元の商品を値引きなし -> 100（割引なし）
   - 100元の商品を0%引き -> 100（定価）
3. 異常系：
   - マイナスの価格 -> エラーを投げること
   - 割引率 > 1 -> エラーを投げること
   - 数値以外の入力 -> エラーを投げること
Vitestフレームワークを使い、完全なテストファイルを示してください。
\end{lstlisting}

この指示の特徴に注目してください。テストケースの3分類（正常·境界·異常）を明示していて、AIが「都合の良いところだけテストする」余地を与えていません。AIはときどき正常系のテストだけ書いて「手早く済ませよう」とします。すべての分類を明示すれば、この問題を避けられます。

\booksection{二、コードレビューのテクニック}

テストでカバーしていても、AIが生成したコードを人間がレビューする作業は欠かせません。コードレビューとは1行ずつ読むことではありません。時間がかかるうえ、あなたのスキルに寄りかかりすぎます。効果的なAIコードレビューでは、次の重要な観点に注目します。

\begin{center}
\begin{tabularx}{\linewidth}{|X|X|X|}
\hline
\textbf{レビュー観点} & \textbf{何に注目するか} & \textbf{よくある問題} \tabularnewline
\hline
セキュリティ & 入力バリデーション、キーの取り扱い、SQLインジェクション対策 & AIは入力チェックを忘れがち。キーをコードにハードコードしてしまう \tabularnewline
エラー処理 & 例外をキャッチしているか、エラーメッセージは適切か & AIは空のcatchブロックやconsole.logだけで済ませがち \tabularnewline
パフォーマンス & N+1クエリや無限ループなど、明らかな性能問題がないか & AIはときどきO(n²)のアルゴリズムを書いてくる \tabularnewline
保守性 & コード構造は明快か、命名は意味を持つか & AIはときどき抽象化しすぎたコードを生成する \tabularnewline
依存関係の管理 & 導入した依存は本当に必要か、バージョンは明確か & AIは不要なライブラリを入れがち \tabularnewline
\hline
\end{tabularx}
\end{center}

\begin{center}\small\color{gray}AIコードレビューの重要な観点\end{center}

レビューですべて自分で1つずつ確認する必要はありません。一次レビューはAIに任せられます。指示の例は「このコードのセキュリティ、エラー処理、パフォーマンスをレビューし、潜在的な問題を指摘してください」。AIのセルフレビュー能力は向上を続けていて、人間のレビューを完全に置き換えることはできませんが、目につく問題の大半をふるい落とすのには役立ちます。

\booksection{三、バージョン管理：AIプログラミングにおけるGitの役割}

AIプログラミングにおいてGitの役割は、従来のプログラミングよりもさらに重要です。AIの変更はどれも予期しない影響を持ちうるので、頼れる「後悔しないための薬」が必要です。バージョン管理こそが、その薬です。

AIプログラミングにおけるGitのベストプラクティスは「こまめにコミット、小刻みに前進」です。検証できる機能がひとつ完成するごとにコミットします。そうすれば、AIのある変更で問題が出ても、直前の使える状態へ正確に戻れ、すべての作業をやり直す必要がありません。

もうひとつの実用的なテクニックが「ブランチでの隔離」です。結果の読めないAIの変更（一括リファクタリングなど）は、まず新しいブランチを作ってそこで行います。結果に満足すればメインブランチへマージし、満足できなければブランチを削除するだけで、メインブランチには影響しません。メインブランチで何度も修正を重ねるより、はるかに安全です。

\begin{enumerate}
  \item AIに修正させる前に、作業ツリーがきれいな状態（未コミットの変更がない）であることを確認します。そうしておけば、AIの修正が気に入らなければ git checkout . ですぐに戻せます。
  \item AIの修正のあと、検証が通ったらすぐにコミットします。コミットメッセージには今回の変更で何をしたかを書きます。たとえば「feat: ユーザー登録機能を追加」のような具合です。
  \item 一括リファクタリングでは、まずブランチを作成します：git checkout -b refactor/auth-system。ブランチ上で修正を完成させ、テストが通ってからマージします。
  \item diffを確認する習慣をつけましょう。AIが修正するたびにgit diffで何が変わったのかを見ます。問題の発見に役立つだけでなく、コードの学びにもなります。
\end{enumerate}

\booksection{四、よくある落とし穴とその回避ガイド}

AIプログラミングには繰り返し現れる落とし穴がいくつかあり、ほぼすべての初心者がそこを踏みます。あらかじめ知っておけば、遠回りをかなり減らせます。

\booksubsection{落とし穴その1：信じすぎの罠}

「AIが書いたコードはプロっぽく見えるから、きっと大丈夫だろう」——これは最も危険な罠です。AIが生成するコードは表面がとてもきれいです。命名は規範的、コメントは揃い、構造は明快。この「表面的なプロっぽさ」が人を油断させます。しかし表面のきれいさは、ロジックの正しさを意味しません。解毒剤：判定は必ずテストによる検証で行い、「見た目が正しそうか」で判断しないこと。

\booksubsection{落とし穴その2：肥大化の罠}

「AIがコードを書いてくれるなら、機能を盛りだくさんにしてもらおう」——これはコードの肥大化を招きます。AIは「おまけ」を付けるのが得意です。ログインページを頼むと、ユーザー管理システム一式を渡してきます。余計な機能は複雑さを増すだけでなく、バグが出る確率も上げます。解毒剤：Simplicity Firstの原則を厳守し、AIに「この機能だけ作って、余計なものは追加しないで」とはっきり伝えること。

\booksubsection{落とし穴その3：理解ズレの罠}

「AIなら私の言いたいことを分かってくれるはず」——そうとは限りません。AIの理解力は確かに高いのですが、曖昧な要件に対してはまだ「推測」します。推測が外れたら、出てくるコードは間違いです。厄介なのは、外れた推測から生まれたコードがたいてい「もっともらく見える」ことで、理解してくれたと錯覚してしまうことです。解毒剤：指示には具体的な詳細を十分に盛り込み、AIが「推測」する余地を減らすこと。

\booksubsection{落とし穴その4：コンテキストの罠}

「前の対話で言ったから、AIは覚えているはず」——おそらく覚えていません。対話が長くなったり新しいセッションを開いたりすると、AIは最初のコンテキストを失います。重要な制約がプロジェクト記憶ファイルに書かれていなければ、その後の変更で無視される可能性が高いのです。解毒剤：重要な制約はプロジェクト記憶ファイルに書き込み、対話の記憶に頼らないこと。

\begin{tcbwarn}{四大落とし穴まとめ}
信じすぎの罠——「正しそう」を鵜呑みにしない。肥大化の罠——AIは「おまけ」を付けたがる。理解ズレの罠——曖昧な要件は「推測」される。コンテキストの罠——会話の記憶は「忘れる」。対応する解毒剤は、テストによる検証、余計なものを削る、細部を十分伝える、記憶ファイルに書き込む、の4つです。

\end{tcbwarn}

\begin{tcbpractice}{やってみよう：プロジェクトに3つの防線を築こう}
すでにプロジェクト1（ツールサイト）に着手しているなら：(1) AIにコア関数の3種類のテストケースを書かせる。(2) 5つの観点（セキュリティ／エラー処理／パフォーマンス／保守性／依存関係）でコードを一度レビューする。(3) gitでプロジェクトを初期化し、最初のコミットをする。3つの防線がそろって、はじめてあなたのプロジェクトは本当の意味であなたのものになります。

\end{tcbpractice}

\begin{tcbsummary}{本章のまとめ}
\begin{itemize}
  \item AIコードの「正しそう」$\neq$「本当に正しい」——表面の品質が高すぎるかえって油断しやすい
  \item AIにテストを書かせる：正常／境界／異常の3種類のケースを明示して「選択的テスト」を防ぐ
  \item レビューの5観点：セキュリティ、エラー処理、パフォーマンス、保守性、依存関係の管理
  \item Gitは「後悔しないための薬」：こまめにコミット、小刻みに前進、ブランチで隔離、diffをこまめに見る
  \item 四大落とし穴と対策：信じすぎ（テストで検証）、肥大化（余計なものは削る）、理解ズレ（細部まで伝える）、コンテキスト（記憶ファイルに書く）
\end{itemize}

\end{tcbsummary}
